\chapter{Deninger's Arithmetic Cohomology as a Chiral Candidate}
\label{v4-arith:w8-a:deninger-chiral-homology}

Deninger's conjectural formalism demands a cohomology theory for
\(\overline{\mathrm{Spec}\,\mathbb{Z}}\) whose regularised Fredholm
determinant equals the completed Riemann zeta function and whose closed
orbits recover the explicit formula. The bar--cobar construction does not
construct this cohomology. It constructs a candidate: the
ordered-incidence homology of \(\mathcal{V}^{\mathrm{prim}}\) (the
Hochschild trace class of the identity on the global arithmetic Iwahori
Hecke category of \(GL_2\)), two formal Tate polar lines, and the
Bost--Connes Euler-prime-side character. Whether this candidate realizes
Deninger's axioms depends on the spectral-flow hypothesis of
\Cref{v4-arith:3d-acs:conj:spectral-flow}. The central result of this
chapter is the obstruction theorem: no unconditional triple
identification exists among Deninger's conjectural cohomology, the
ordered-incidence candidate, and the boundary homology of the
arithmetic Chern--Simons bulk.

\section{Deninger's Wish List}
\label{v4-arith:w8-a:deninger-wishlist}

Deninger's formalism asks for a cohomology theory
\[
H^\bullet_\Delta(\overline{\mathrm{Spec}\,\mathbb Z})
\]
with three cohomological degrees, a Frobenius/flow operator \(\Theta\)
whose regularised determinant has the non-trivial zero divisor, a
Lefschetz trace formula whose closed-orbit side is the prime side of
the explicit formula, Poincar\'{e} duality, and purity. Under
zero-placement this may be written as
\(\Theta=\tfrac12+iH\) with \(H\) self-adjoint. For the Riemann zeta
function the middle degree should carry the non-trivial zero divisor.

In the next display \(\zeta(s)\) follows Deninger's convention: it is
the completed zeta function, not the bare Euler product.
\begin{equation}
\label{v4-arith:w8-a:eq:deninger-det}
\zeta(s)
\stackrel{?}{=}
\det\nolimits_\infty
\bigl(s\cdot\mathrm{id}-\Theta\mid
H^\bullet_\Delta\bigr)^{(-1)^\bullet}.
\end{equation}

Here \(\det_\infty\) is Deninger's regularised Fredholm determinant,
\(\zeta(s)\) denotes the completed Hasse--Weil zeta (including the
archimedean factor), and the sign \((-1)^\bullet\) is the
Euler-characteristic grading. The
Lefschetz trace formula takes the form
\begin{equation}
\label{v4-arith:w8-a:eq:deninger-lefschetz}
\sum_\rho h(\rho)
\stackrel{?}{=}
\sum_\bullet(-1)^\bullet
\operatorname{tr}\bigl(h(\Theta)\mid H^\bullet_\Delta\bigr)
\stackrel{?}{=}
\text{polar terms}+\text{archimedean term}
-\sum_{p,k}(\log p)p^{-k/2}\widehat{h}(k\log p),
\end{equation}
where \(h\) is a Schwartz test function on \(\mathbb{R}\) and
\(\widehat{h}\) its Fourier transform.

\begin{remark}[absence of a construction]
\label{v4-arith:w8-a:deninger-wishlist:rem:absence}
The axioms above are not a construction. They specify what the missing
arithmetic cohomology must do. The construction provides the prime-side
Euler data (via Bost--Connes) and the formal polar lines (by
insertion); the zero-side operator is not constructed.
\end{remark}

\begin{remark}[Vol IV scoping]
\label{v4-arith:w8-a:deninger-wishlist:rem:vol4}
Chapter~\ref{v4-arith:3d-acs-E1bar} isolates the correct boundary:
the ordered-incidence bar is theorem-grade as an algebraic model, while
the comparison with genuine BD/CG chiral homology, the profinite
quantum ACS theory, and the Deninger spectral flow remain conditional.
This chapter keeps that boundary fixed.
\end{remark}

\section{Surviving Inputs}
\label{v4-arith:w8-a:machinery}
\label{v4-arith:w8-a:machinery:ch23}

The bar--cobar construction produces four inputs with theorem-grade form;
the remainder of the Deninger axioms are either conditional or open.

\begin{enumerate}
\item The finite-level arithmetic Chern--Simons state sum and gluing
formula of \Cref{v4-arith:3d-acs:thm:finite-level-ACS}. This is a
closed finite groupoid state-sum statement, not a full profinite TQFT.
\item The ordered-incidence bar
\Cref{v4-arith:3d-acs:defn:E1-bar} and its normalized-bar computation
\Cref{v4-arith:3d-acs:thm:chi-hom-bar}. This computes the explicitly
defined incidence model, not automatically BD/CG factorization
homology on \(\overline{\mathrm{Spec}\,\mathbb Z}\).
\item The Bost--Connes Euler skeleton
\Cref{v4-arith:3d-acs:prop:deninger-prime-side}. It gives the KMS
character \(\zeta(s)\) and the prime-power logarithmic derivative in
\(\mathrm{Re}(s)>1\).
\item The Deninger spectral-flow hypothesis
\Cref{v4-arith:3d-acs:conj:spectral-flow}: the zero-side determinant
divisor, trace distribution, and Hilbert--P\'olya self-adjoint
determinant-trace operator. This is the single missing piece that the
bar--cobar construction does not supply.
\end{enumerate}

\section{The Chiral Candidate}
\label{v4-arith:w8-a:identification}

Given the four inputs above, one may assemble a graded object that
mimics Deninger's three-degree shape. The question is whether this
assembly is genuinely Deninger's cohomology or only a formal
approximation to it.

Let \(B^{\mathrm{ord}}_\bullet(\mathcal{V}^{\mathrm{prim}})\) denote the
ordered-incidence bar complex of \Cref{v4-arith:3d-acs:defn:E1-bar}
applied to \(\mathcal{V}^{\mathrm{prim}}\).

\begin{definition}[arithmetic chiral homology candidate]
\label{v4-arith:w8-a:def:chiral-homology}
\ClaimStatusDefinitional. The \emph{arithmetic chiral homology
candidate} is
\begin{equation}
\label{v4-arith:w8-a:eq:chiral-homology-def}
H^\bullet_{\mathrm{ord}}
(\overline{\mathrm{Spec}\,\mathbb Z},\mathcal{V}^{\mathrm{prim}})
:=
H^\bullet\bigl(
B^{\mathrm{ord}}_\bullet(\mathcal V^{\mathrm{prim}})
\bigr).
\end{equation}
It is an ordered-incidence model. Calling it genuine arithmetic chiral
homology requires the separate BD/CG comparison of
\Cref{v4-arith:3d-acs:cor:BDCG-ordered-comparison}.
\end{definition}

\begin{definition}[Deninger chiral candidate]
\label{v4-arith:w8-a:def:deninger-cohomology}
\ClaimStatusDefinitional. The \emph{Deninger chiral candidate} is the
candidate complex of \Cref{v4-arith:3d-acs:def:chiral-deninger-complex}:
\begin{equation}
\label{v4-arith:w8-a:eq:deninger-chiral-def}
H^\bullet_{\Delta,\chi}
(\overline{\mathrm{Spec}\,\mathbb Z})
:=
H^\bullet_{\mathrm{Den},\chi}
(\overline{\mathrm{Spec}\,\mathbb Z}).
\end{equation}
It consists of the ordered-incidence core, completed continuous duals,
and formal Tate polar lines in degrees \(0\) and \(2\).
\end{definition}

\begin{remark}[three labels, not one theorem]
\label{v4-arith:w8-a:rem:three-labels}
Three nearby objects appear: Deninger's conjectural cohomology,
the ordered-incidence chiral candidate, and the boundary state space of
the conjectural 3D ACS bulk. Their equality is a conditional assertion,
not a proved identification.
\end{remark}

\begin{proposition}[conditional Deninger comparison]
\label{v4-arith:w8-a:thm:deninger-chiral-equals}
\ClaimStatusProvedHereConditional \emph{(on
\Cref{v4-arith:3d-acs:conj:spectral-flow}, the BD/CG comparison
\Cref{v4-arith:3d-acs:cor:BDCG-ordered-comparison}, and the
topological-dual choices in
\Cref{v4-arith:3d-acs:def:chiral-deninger-complex})}. The chiral
candidate \(H^\bullet_{\Delta,\chi}\), together with the two Tate polar
lines and the spectral-flow operator assumed in
\Cref{v4-arith:3d-acs:conj:spectral-flow}, carries Deninger's formal
axioms for \(\overline{\mathrm{Spec}\,\mathbb Z}\).
\end{proposition}

\begin{proof}
The finite-prime Euler terms are exactly the Bost--Connes series of
\Cref{v4-arith:3d-acs:prop:deninger-prime-side}. The polar lines are
inserted by definition in
\Cref{v4-arith:3d-acs:def:chiral-deninger-complex}. The zero-side
operator, determinant identity, trace formula, and purity are precisely
the extra data demanded by
\Cref{v4-arith:3d-acs:conj:spectral-flow}. Thus the statement is a
conditional implication, not an independent construction of the missing
cohomology.
\end{proof}

\section{No Triple Identification Yet}
\label{v4-arith:w8-a:integration}
\label{v4-arith:w8-a:ah:identification}

\begin{theorem}[no unconditional triple identification]
\label{v4-arith:w8-a:thm:triple-id}
\ClaimStatusProvedHere. No unconditional equality holds among
Deninger's conjectural cohomology, the ordered-incidence chiral
candidate, and the boundary homology of \(\mathrm{ACS}^{\mathrm{qu}}\).
\end{theorem}

\begin{proof}
The first object is conjectural. The second is an explicitly defined
ordered-incidence homology object. The third depends on the profinite
quantum ACS conjecture and on a boundary comparison with
\(\mathcal V^{\mathrm{prim}}\). The only theorem-grade equality is the
normalized-bar computation of the ordered-incidence model. The BD/CG
comparison and spectral-flow identifications are listed as residual
conditions in \Cref{v4-arith:3d-acs:residual}.
\end{proof}

\begin{corollary}[conditional triangle]
\label{v4-arith:w8-a:cor:triangle}
\ClaimStatusConditional \emph{(on the profinite ACS conjecture, the
BD/CG ordered-model comparison, a regularised character map, and the
Deninger spectral-flow hypothesis)}. These assumptions yield
\begin{equation}
\label{v4-arith:w8-a:eq:triangle}
Z_{\mathrm{ACS}^{\mathrm{qu}}}
(\overline{\mathrm{Spec}\,\mathbb Z};k^{\mathrm{Ar}})
=
\xi(k^{\mathrm{Ar}})
=
\operatorname{ch}_{\mathrm{reg}}
\bigl(H^\bullet_{\Delta,\chi}\bigr)(k^{\mathrm{Ar}}),
\end{equation}
where \(\xi(s) = \tfrac{1}{2}s(s-1)\pi^{-s/2}\Gamma(s/2)\zeta(s)\)
is the completed Riemann zeta function satisfying \(\xi(s)=\xi(1-s)\),
\(k^{\mathrm{Ar}}\) is the arithmetic level parameter, and
\(\operatorname{ch}_{\mathrm{reg}}\) denotes the regularised graded
trace (the \(\Theta\)-eigenvalue generating series with Deninger
regularisation). Without those assumptions,
\eqref{v4-arith:w8-a:eq:triangle} is a conditional shorthand.
\end{corollary}

\section{Frobenius and Flow}
\label{v4-arith:w8-a:global-F}

The Bost--Connes Euler skeleton provides prime-length data at finite
primes. Deninger's formalism requires a global Frobenius/flow operator
acting on all three cohomological degrees. The question is whether the
Euler skeleton and the archimedean scaling generator combine into such
an operator.

\begin{definition}[finite Euler operators]
\label{v4-arith:w8-a:def:F-finite}
\ClaimStatusDefinitional. At a finite prime \(p\), the Bost--Connes
Adams operation \(\psi_p\) records the prime length \(\log p\) and
contributes the Euler factor \((1-p^{-s})^{-1}\). It is denoted here by
\[
F^{\mathrm{Euler}}_p:=\psi_p.
\]
This is an Euler-skeleton operator, not by itself Deninger's global
Frobenius on middle cohomology.
\end{definition}

\begin{proposition}[finite Wilson/Frobenius comparison is conditional]
\label{v4-arith:w8-a:lem:wilson-finite}
\ClaimStatusConditional \emph{(on the Wilson-line axiom in
\Cref{v4-arith:3d-acs:conj:qu-acs})}. A Wilson loop around the
Morishita knot attached to \(p\) should recover the corresponding
Frobenius trace. This is part of the conjectural ACS datum, not a
consequence of the ordered bar.
\end{proposition}

\begin{definition}[archimedean scaling candidate]
\label{v4-arith:w8-a:def:F-infinity}
\ClaimStatusDefinitional. The archimedean scaling generator is the
Mellin infinitesimal \(-\partial_t\) on \(\mathbb{R}_+\). The
candidate archimedean flow on \(H^\bullet_{\Delta,\chi}\) is denoted
\(\Theta_{\infty,\chi}\). In Deninger's and Connes' formulation, the
target operator has the form
\(\Theta_\xi = \tfrac{1}{2}\mathrm{id} + iH_{\mathrm{HP}}\),
where \(H_{\mathrm{HP}}\) is a self-adjoint Hilbert--Polya operator
(to be constructed) whose normal determinant divisor is
\(\mathrm{div}\,\xi_{\mathbb R}\). After zero-placement its
Weil--Connes trace is supported at the ordinate multiset
\(\{\gamma_\rho\}\). The identification
\begin{equation}
\label{v4-arith:w8-a:eq:F-infinity}
\Theta_{\infty,\chi}\stackrel{?}{=}
\Theta_\xi
\end{equation}
is the spectral-flow hypothesis and is not available unconditionally.
\end{definition}

\begin{proposition}[Virasoro conjugacy is not available]
\label{v4-arith:w8-a:lem:L0-conjugacy}
\ClaimStatusConditional \emph{(on the spectral-flow hypothesis)}.
Without the spectral-flow theorem of
\Cref{v4-arith:3d-acs:conj:spectral-flow}, the Virasoro
\(L_0^{\mathrm{Ar}}\), the Bost--Connes logarithmic Hamiltonian, and
the Hilbert--Polya operator \(H_{\mathrm{HP}}\) remain distinct
operators. Their identification requires the spectral-flow theorem as
input.
\end{proposition}

\begin{definition}[formal global flow]
\label{v4-arith:w8-a:def:F-global}
\ClaimStatusDefinitional. The formal global flow symbol is
\begin{equation}
\label{v4-arith:w8-a:eq:F-global}
F_\chi:=
\Theta_{\infty,\chi}\oplus\bigoplus_p F^{\mathrm{Euler}}_p.
\end{equation}
The notation records finite-prime lengths and archimedean scaling. It
does not assert a well-defined Deninger Frobenius on a completed
cohomology theory.
\end{definition}

\begin{theorem}[global Frobenius obstruction]
\label{v4-arith:w8-a:thm:F-global-welldef}
\ClaimStatusProvedHere. The ordered-incidence bar and Bost--Connes
Euler skeleton do not construct Deninger's global Frobenius operator.
\end{theorem}

\begin{proof}
The Bost--Connes side gives Euler factors and the prime-power series
in the half-plane of convergence. Deninger's operator must act on the
middle cohomology, must have regularised determinant divisor
\(\mathrm{div}\,\xi_{\mathbb R}\), and must participate in a
Lefschetz trace formula. That zero-side operator is exactly the missing
spectral-flow input of
\Cref{v4-arith:3d-acs:conj:spectral-flow}.
\end{proof}

\begin{remark}[archimedean Wilson surface]
\label{v4-arith:w8-a:rem:wilson-infinity}
The phrase ``archimedean Wilson surface'' is useful only as a
conditional dictionary entry inside the profinite ACS conjecture. It
is not an independent construction of \(\Theta_\xi\).
\end{remark}

\section{Character, Degrees, and Purity}
\label{v4-arith:w8-a:character}
\label{v4-arith:w8-a:H-structure}

The Frobenius obstruction theorem leaves open the question of what the
candidate complex does compute. The regularised graded trace and the
formal Tate structure answer this question partially: the Euler product
and the polar lines are present; the zero-side is missing.

\begin{theorem}[character identity is conditional]
\label{v4-arith:w8-a:thm:character-identity}
\ClaimStatusConditional \emph{(on
\Cref{v4-arith:3d-acs:cor:partition-bar} and
\Cref{v4-arith:3d-acs:conj:spectral-flow})}. The regularised
character of the Deninger chiral candidate is \(\xi(s)\) only after
the partition-function comparison and the zero-side spectral-flow
operator are assumed.
\begin{equation}
\label{v4-arith:w8-a:eq:character}
\operatorname{ch}_{\mathrm{reg}}
\bigl(H^\bullet_{\Delta,\chi}\bigr)(s)
\stackrel{\mathrm{cond}}{=}
\xi(s).
\end{equation}
\end{theorem}

\begin{remark}[normalisation]
\label{v4-arith:w8-a:rem:normalisation}
The additive determinant form
\(\det_\infty(s\cdot\mathrm{id}-\Theta)\) and the multiplicative
Euler-character form are interchangeable only after the operator
\(\Theta\), the polar signs, and the regularisation convention have
been fixed.
\end{remark}

\begin{proposition}[formal Tate polar lines]
\label{v4-arith:w8-a:prop:H0-H2}
\ClaimStatusProvedHere. The candidate complex contains formal
one-dimensional Tate lines in degrees \(0\) and \(2\). This accounts
for the polar terms in the Deninger skeleton but does not prove that
the total complex has only Deninger's three degrees.
\end{proposition}

\begin{proposition}[middle cohomology is the hard part]
\label{v4-arith:w8-a:prop:H1-structure}
\ClaimStatusConditional \emph{(on
\Cref{v4-arith:3d-acs:conj:spectral-flow})}. The middle cohomology
\(H^1_{\Delta,\chi}\) admits a Hilbert-space completion with a
self-adjoint generator having spectrum \(\{\gamma_\rho\}\) only
after \Cref{v4-arith:3d-acs:conj:spectral-flow} is assumed. This is
the Hilbert--Polya hypothesis in cohomological form.
\end{proposition}

\begin{remark}[critical-line locus]
\label{v4-arith:w8-a:rem:crit-line-weight}
Purity of the middle degree is the assertion that the zero-side
operator has real part \(1/2\). For zeta this is RH in cohomological
language.
\end{remark}

\section{Duality and Positivity}
\label{v4-arith:w8-a:poincare-duality}

Koszul duality for the arithmetic bar complex gives the
functional-equation symmetry \(s\leftrightarrow 1-s\). The question is
whether this duality implies Poincar\'{e} duality for the Deninger
cohomology, and specifically whether it forces the middle degree to be
positive definite.

\begin{theorem}[duality gives the functional equation, not positivity]
\label{v4-arith:w8-a:thm:triple-duality}
\ClaimStatusConditional \emph{(on the relevant Koszul and ACS
comparison hypotheses)}. The Koszul self-duality framework accounts
for the functional-equation symmetry \(s\leftrightarrow 1-s\). It does
not by itself imply positivity of the Deninger middle pairing or RH.
\end{theorem}

\section{Purity as the RH-Level Input}
\label{v4-arith:w8-a:purity}

Duality accounts for the functional equation but not for the
zero-placement. Purity of the middle cohomological degree is the
assertion that the zero-side operator \(\Theta_\xi\) is self-adjoint
with spectrum on the line \(\mathrm{Re}(s)=1/2\), which is RH.

\begin{theorem}[purity is RH-level]
\label{v4-arith:w8-a:thm:purity-equivalence}
\ClaimStatusConditional \emph{(on the spectral-flow/Hilbert--Polya
realisation)}. Deninger purity for
\(H^1_{\Delta,\chi}\), Hilbert--Polya self-adjointness for
\(\Theta_\xi\), and the RH zero-placement statement are equivalent
forms of the same missing input.
\end{theorem}

\begin{corollary}[Deninger-chiral RH construction]
\label{v4-arith:w8-a:cor:rh-construction}
\ClaimStatusConditional. A proof of RH by this route is precisely a
proof of the spectral-flow/positivity input. The ordered-incidence bar
and the Bost--Connes Euler skeleton reduce the shape of the problem;
they do not remove it.
\end{corollary}

\section{All-Genus Extension}
\label{v4-arith:w8-a:all-genus}

The preceding sections address
\(\overline{\mathrm{Spec}\,\mathbb{Z}}\) and the Riemann zeta function.
In the notation of
\Cref{v4-arith:ag:defn:datum} this is the zeta fibre of the
rank-graded arithmetic family, not a smooth surface of curve genus
zero. Higher-rank arithmetic data occur in several distinct windows:
arithmetic surfaces, totally imaginary number fields, crystalline
deformation stacks, Hodge--Arakelov theta data, and Shimura varieties.
The formal extension records only what the ordered candidate produces
after such a datum has been fixed.

\begin{definition}[genus-\(g\) chiral Deninger candidate]
\label{v4-arith:w8-a:def:deninger-g}
\ClaimStatusDefinitional. A \emph{genus-\(g\) arithmetic datum}
\(\mathcal{D}_g\) means the datum of
\Cref{v4-arith:ag:defn:datum}. Thus it may be an arithmetic surface
only in Windows~I and~V; in the other windows it is respectively a
number-field, crystalline, Hodge--Arakelov, or Shimura datum. The
corresponding \(\mathcal{V}^{\mathrm{prim},g}\) is the row-specific
Hochschild trace class of the identity on the indicated arithmetic
Hecke category. In the surface, crystalline, and automorphic rows its
oscillator sector has rank \(2g\); in the Dedekind row \(g\) records
the complex embedding rank of the coefficient field. The
\emph{genus-\(g\) chiral Deninger candidate} is
\begin{equation}
\label{v4-arith:w8-a:eq:deninger-g-def}
H^\bullet_{\Delta,\chi,g}
:=
H^\bullet_{\mathrm{ord}}
(\mathcal D_g,\mathcal V^{\mathrm{prim},g}).
\end{equation}
The weight of \(\mathcal{D}_g\), denoted \(w_{\mathcal{D}_g}\), is the
motivic weight of its \(L\)-function; purity is the assertion that
\(\mathrm{Re}(s)=w_{\mathcal{D}_g}/2\) for all non-trivial zeros.
\end{definition}

\begin{theorem}[all-genus Deninger comparison]
\label{v4-arith:w8-a:thm:all-genus-deninger}
\ClaimStatusConditional \emph{(on the all-genus ACS, ordered-model
comparison, and spectral-flow hypotheses)}. The genus-\(g\) candidate
carries the expected motivic functional-equation centre and prime-side
Euler data. It does not construct the genus-\(g\) Deninger cohomology
without the same spectral-flow input.
\end{theorem}

\begin{remark}[weight convention]
\label{v4-arith:w8-a:rem:weight-g}
The expected purity line is
\(\mathrm{Re}(s)=w_{\mathcal D_g}/2\). This is a target condition, not
a theorem of the ordered bar.
\end{remark}

\section{Dictionary}
\label{v4-arith:w8-a:dictionary}

Table~\ref{v4-arith:w8-a:table:triple} records which Deninger axioms
the chiral candidate supplies and which remain open or conditional.

\begin{table}[h]
\centering
\small
\begin{tabular}{p{0.26\linewidth}p{0.31\linewidth}p{0.31\linewidth}}
\toprule
Deninger target & Chiral/ACS object & Form \\
\midrule
finite prime orbit lengths & Bost--Connes Euler skeleton & proved in
\(\mathrm{Re}(s)>1\) \\
cohomology object & ordered-incidence candidate plus polar lines &
defined, not identified with Deninger cohomology \\
global Frobenius/flow & spectral-flow operator \(\Theta_\xi\) &
open, Hilbert--Polya strength \\
determinant \(=\xi\) & regularised character & conditional \\
Lefschetz trace formula & prime-power trace distribution & open \\
purity & real part \(1/2\) for middle spectrum & RH-level \\
\bottomrule
\end{tabular}
\caption{Surviving Deninger--chiral dictionary.}
\label{v4-arith:w8-a:table:triple}
\end{table}

\section{Unconditional Deninger--Chiral Consequences}
\label{v4-arith:w8-a:what-closes}

The ordered-incidence construction proves the four assertions below;
the six remaining comparison maps require the stated hypotheses.

\subsection*{Unconditional consequences}
\label{v4-arith:w8-a:what-closes:uncond}

\begin{enumerate}
\item Ordered-incidence homology is computed by the normalized ordered
bar.
\item The finite-prime Bost--Connes side gives the Euler product and
logarithmic derivative in the convergence half-plane.
\item The formal Tate polar lines can be inserted as part of the
candidate complex.
\item Koszul self-duality explains the functional-equation symmetry.
\end{enumerate}

\subsection*{Conditional}
\label{v4-arith:w8-a:what-closes:cond}

\begin{enumerate}
\item Profinite quantum arithmetic Chern--Simons.
\item Boundary comparison with \(\mathcal V^{\mathrm{prim}}\).
\item BD/CG comparison with the ordered-incidence model.
\item Regularised character equality with \(\xi\).
\item Spectral flow/Hilbert--P\'olya zero-divisor operator.
\item Positivity/purity of the middle degree.
\end{enumerate}

\section{Residual Conditions}
\label{v4-arith:w8-a:residual}

The residue is a single mathematical obstruction: construct the middle
operator and prove the trace formula with positivity. Equivalently,
prove the spectral-flow hypothesis of
\Cref{v4-arith:3d-acs:conj:spectral-flow}. The remaining assertions are
unconditional consequences or conditional consequences of that
hypothesis.
